% Auto-generated complete chapter from Research/cerrando_sistema.md
\chapter{Estado actual}
\label{complete:research_051}

\begin{sourcemeta}
\textbf{Fuente:} \href{file:///E:/Laboral/Research/cerrando_sistema.md}{Research/cerrando\_sistema.md}\\
\textbf{Seccion:} research\\
\textbf{Estado:} research\_snapshot\\
\textbf{Rol:} Research / estudio de mercado\\
\textbf{Lineas originales:} 340\\
\textbf{Ultima modificacion:} 2026-06-11 15:11\\
\textbf{Deduplicacion conservadora:} 0 bloques repetidos omitidos; 0 caracteres repetidos no reimpresos.
\end{sourcemeta}

\section*{Resumen de orientacion}
Sí. Falta \textbf{cerrar el sistema}, no necesariamente investigar mucho más. Ya tienes una base modular suficientemente amplia, pero todavía hay brechas entre \textbf{investigación} y \textbf{ejecución real}.

\section*{Contenido completo depurado}
\addcontentsline{toc}{section}{Contenido completo depurado}

Sí. Falta \textbf{cerrar el sistema}, no necesariamente investigar mucho más. Ya tienes una base modular suficientemente amplia, pero todavía hay brechas entre \textbf{investigación} y \textbf{ejecución real}.




\subsection{Estado actual}




Ya está cubierto el bloque estratégico principal:




\begin{mdcode}
Bloque de codigo: text
00_master_audit_blueprint.md
01_source_of_truth_profile.md
02_market_role_fit_and_positioning.md
03_salary_benchmark_and_remote_colombia.md
04_eu_portugal_germany_mobility_strategy.md
05_language_strategy_and_roi.md
06_education_courses_rebelway_masters_phd.md
07_portfolio_strategy_and_project_architecture.md
08_twinsight_x500_case_study.md
09_ara_ai_tools_and_secondary_projects.md
10_github_linkedin_cv_strategy.md
11_company_targets_job_boards_recruiters.md
12_application_system_and_tracker.md
13_interview_negotiation_and_offer_risk.md
14_30_60_90_execution_plan.md
15_final_integrated_strategy.md
\end{mdcode}




Eso ya funciona como \textbf{investigación base + plan estratégico}.




Pero todavía no es un sistema completamente ejecutable porque faltan artefactos operativos.




\medskip\hrule\medskip




\section{Qué falta}




\subsection{1. Auditoría final de consistencia}




Antes de producir CV, LinkedIn, portfolio y trackers, conviene hacer una pasada final:




\begin{mdcode}
Bloque de codigo: text
16_final_consistency_audit_and_gap_register.md
\end{mdcode}




Debe revisar:




\begin{itemize}
\item contradicciones entre secciones;
\item claims peligrosos;
\item datos que quedaron débiles;
\item secciones que dependen de verificación futura;
\item qué archivos son estrategia y cuáles son ejecución;
\item qué acciones deben hacerse antes de aplicar;
\item qué cosas no se deben afirmar públicamente.
\end{itemize}




\textbf{Herramientas:} no Deep Research, no agente.




Esto sería el cierre lógico antes de pasar a producción de materiales.




\medskip\hrule\medskip




\subsection{2. Archivos ejecutables que todavía faltan}




La investigación dice qué hacer, pero faltan documentos de uso directo:




\begin{mdcode}
Bloque de codigo: text
17_cv_base_and_role_variants.md
18_linkedin_final_rewrite.md
19_github_readme_templates.md
20_portfolio_copy_and_site_structure.md
21_twinsight_demo_video_script_and_shotlist.md
22_outreach_templates_english_spanish.md
23_interview_answer_bank.md
24_offer_evaluation_scorecard.md
25_application_tracker_template.csv / .xlsx
26_weekly_execution_dashboard.md
\end{mdcode}




Estos son los que realmente convierten la estrategia en acción.




\medskip\hrule\medskip




\subsection{3. Validación externa pendiente}




La sección con más incertidumbre sigue siendo:




\begin{mdcode}
Bloque de codigo: text
11_company_targets_job_boards_recruiters.md
\end{mdcode}




Ahí sí hay riesgo de que algunos datos estén incompletos o no actualizados, porque empresas, careers pages y remote policies cambian constantemente.




No intentaría perfeccionar las 120 empresas ahora. Haría esto:




\subsubsection{Pasada práctica}




Elegir solo:




\begin{mdcode}
Bloque de codigo: text
Top 30 empresas A-priority
Top 20 empresas B-priority
Top 20 job boards / recruiters
\end{mdcode}




Y verificar:




\begin{itemize}
\item careers URL;
\item si hay roles vivos;
\item si aceptan remoto;
\item si remoto significa worldwide, LATAM, US-only o EU-only;
\item si hay contractor/B2B;
\item idioma;
\item seniority;
\item si TwinSight encaja.
\end{itemize}




\textbf{Herramienta recomendada:} modo agente, pero solo para esa verificación concreta.




No usaría agente para todo el documento.




\medskip\hrule\medskip




\section{Qué haría ahora}




\subsection{Paso 1 — Generar \texttt{16\_final\_consistency\_audit\_and\_gap\_register.md}}




Este archivo debe responder:




\begin{mdcode}
Bloque de codigo: text
¿Qué está cerrado?
¿Qué falta?
¿Qué tiene riesgo?
¿Qué requiere verificación humana?
¿Qué requiere agente?
¿Qué debe producirse como material final?
¿Qué se puede ejecutar ya?
¿Qué no se debe hacer todavía?
\end{mdcode}




Esto crea un puente entre la investigación y los materiales finales.




\textbf{Herramientas:} ninguna.




\medskip\hrule\medskip




\subsection{Paso 2 — Crear el paquete de aplicación}




Después de aprobar el gap register, generaría en este orden:




\subsubsection{1. CV}




\begin{mdcode}
Bloque de codigo: text
17_cv_base_and_role_variants.md
\end{mdcode}




Debe incluir:




\begin{itemize}
\item CV base;
\item CV para Real-Time 3D / Unity;
\item CV para Technical Visualization / Digital Twin;
\item CV para Unity WebGL / Interactive 3D;
\item CV para Tools / Python Automation;
\item bullet bank;
\item reglas de edición.
\end{itemize}




\subsubsection{2. LinkedIn}




\begin{mdcode}
Bloque de codigo: text
18_linkedin_final_rewrite.md
\end{mdcode}




Debe incluir:




\begin{itemize}
\item headline final;
\item About;
\item Experience;
\item Education;
\item Projects;
\item Featured;
\item Skills;
\item recruiter keywords;
\item availability statement.
\end{itemize}




\subsubsection{3. GitHub}




\begin{mdcode}
Bloque de codigo: text
19_github_readme_templates.md
\end{mdcode}




Debe incluir:




\begin{itemize}
\item README TwinSight;
\item README ARA;
\item repo descriptions;
\item pinned repo order;
\item disclaimer de IA;
\item disclaimer de assets/CAD;
\item GitHub profile README.
\end{itemize}




\subsubsection{4. Portfolio}




\begin{mdcode}
Bloque de codigo: text
20_portfolio_copy_and_site_structure.md
\end{mdcode}




Debe incluir:




\begin{itemize}
\item copy listo para web;
\item home page;
\item about;
\item project cards;
\item case study TwinSight;
\item ARA;
\item Blender portrait;
\item contact;
\item estructura MVP.
\end{itemize}




\subsubsection{5. Demo video}




\begin{mdcode}
Bloque de codigo: text
21_twinsight_demo_video_script_and_shotlist.md
\end{mdcode}




Debe incluir:




\begin{itemize}
\item guion de 60 s;
\item guion de 120 s;
\item lista de tomas;
\item orden visual;
\item captions;
\item versión sin voz;
\item versión para LinkedIn.
\end{itemize}




\medskip\hrule\medskip




\section{Lo que falta como información tuya}




Para cerrar con máxima precisión, faltan datos concretos:




\subsection{TwinSight}




\begin{itemize}
\item URL final del repo público.
\item URL final del demo WebGL si existe.
\item Capturas finales.
\item Video o grabación.
\item Métricas finales verificadas.
\item Texto exacto aprobado de la tesis.
\item Qué partes del CAD pueden mostrarse públicamente.
\item Qué assets deben tener disclaimer.
\end{itemize}




\subsection{CV / LinkedIn}




\begin{itemize}
\item Historial laboral exacto con fechas.
\item Nombre de clientes freelance si pueden mencionarse.
\item Si hay proyectos pagados que puedan mostrarse.
\item Nivel real de inglés: certificado o solo autodeclarado.
\item Fecha estimada de sustentación.
\item Disponibilidad semanal exacta después de sustentación.
\end{itemize}




\subsection{Legal / contractor}




\begin{itemize}
\item Si ya tienes RUT.
\item Si ya facturas electrónicamente.
\item Si tienes contador.
\item Método preferido de pago.
\item Si aceptarías contrato por hora, mensual retainer o fixed-price.
\item Mínimo neto real aceptable.
\end{itemize}




\subsection{Portafolio}




\begin{itemize}
\item Dominio o plataforma elegida.
\item Si quieres GitHub Pages, Framer, Webflow, Notion, HTML propio o WordPress.
\item Estilo visual preferido.
\item Si quieres portfolio en inglés solamente o bilingüe.
\end{itemize}




\medskip\hrule\medskip




\section{Roadmap seguro}




\subsection{Fase A — Cierre documental}




\begin{mdcode}
Bloque de codigo: text
16_final_consistency_audit_and_gap_register.md
\end{mdcode}




Objetivo: asegurar que nada se contradice.




\subsection{Fase B — Materiales públicos}




\begin{mdcode}
Bloque de codigo: text
17_cv_base_and_role_variants.md
18_linkedin_final_rewrite.md
19_github_readme_templates.md
20_portfolio_copy_and_site_structure.md
21_twinsight_demo_video_script_and_shotlist.md
\end{mdcode}




Objetivo: quedar listo para aplicar.




\subsection{Fase C — Sistema de aplicación}




\begin{mdcode}
Bloque de codigo: text
22_outreach_templates_english_spanish.md
23_interview_answer_bank.md
24_offer_evaluation_scorecard.md
25_application_tracker_template.xlsx
26_weekly_execution_dashboard.md
\end{mdcode}




Objetivo: operar la búsqueda laboral.




\subsection{Fase D — Verificación externa puntual}




Usar modo agente para:




\begin{mdcode}
Bloque de codigo: text
Top 30 companies verification
Top 20 job boards verification
Top recruiter list verification
\end{mdcode}




Objetivo: no aplicar a listas muertas o mal filtradas.




\subsection{Fase E — Ejecución 30 días}




Aplicar, medir, corregir.




\medskip\hrule\medskip




\section{Recomendación inmediata}




El siguiente archivo debería ser:




\begin{mdcode}
Bloque de codigo: text
16_final_consistency_audit_and_gap_register.md
\end{mdcode}




\textbf{Herramientas:}




\begin{center}
\scriptsize
\begin{longtable}{>{\raggedright\arraybackslash}p{0.455\linewidth} >{\raggedright\arraybackslash}p{0.455\linewidth}}
\toprule
Herramienta & Usar \\
\midrule
\endfirsthead
\toprule
Herramienta & Usar \\
\midrule
\endhead
Deep Research & No \\
Modo agente & No \\
\bottomrule
\end{longtable}
\end{center}




Después de ese archivo, pasamos a generar materiales finales de aplicación. Ahí ya no estaremos investigando: estaremos produciendo el kit real para buscar trabajo.



